\documentclass[12pt]{nextart_zh}
\UseFont{tibetan}\UseFont{libertinus}\UseFeature{hyperlinks}\UseFeature{tables}
\title{藏文轉寫與 GHC 對照簡表}\author{Kelvin Yang}\date{2026年8月}
\begin{document}\maketitle
本表將本方案使用的藏文結構與龔煌城擬音（GHC）對照。藏文拼寫保存聲韻地位；GX 與 GHC 則對同一批正字特徵作不同的音系分析。

\section{元音與等第}
\begin{center}\begin{tabular}{lll}\toprule 藏文 & GX & GHC \\\midrule
ཨ & \emph{Ca} & \emph{Cja} \\ ཨེ & \emph{Ce} & \emph{Cjij} \\ ཨི & \emph{Ci} & \emph{Cji} \\ ཨོ & \emph{Co} & \emph{Cjo} \\ ཨུ & \emph{Cu} & \emph{Cju} \\ ཨྀ & \emph{Cə} & \emph{Cjɨ} \\
ཨའ & \emph{Ca̱} & \emph{Ca} \\ ཨའེ & \emph{Ce̱} & \emph{Cej} \\ ཨའི & \emph{Ci̱} & \emph{Ce} \\ ཨའོ & \emph{Co̱} & \emph{Co} \\ ཨའུ & \emph{Cu̱} & \emph{Cu} \\ ཨའྀ & \emph{Cə̱} & \emph{Cə} \\\bottomrule
\end{tabular}\end{center}

\section{主要結構}
\begin{center}\small\begin{tabular}{lll}\toprule 藏文手段 & 聲韻地位 & GHC 結果 \\\midrule
後置འ & 一等 & 無 \emph{j-} 的元音系列 \\ 後置ཡ & 四等 & 在 GX/GHC 與三等合流 \\ 前置འ- & marked/column-2 & 相應長元音系列 \\ ར- & 捲舌韻 & 韻尾 \emph{-r} \\ 上加ས- & 緊韻 & 元音下點 \\ -མ/-ག༹ & GX \emph{-ṃ/-w} & 鼻化、\emph{-j/-w} 韻 \\ ཝ；stack 中ྭ & GX \emph{v/vw} & GHC \emph{w}；條件性 \emph{w} 不重寫 \\ \texttt{?}/\texttt{†} & 聲調存疑/其他存疑 & 分別保留 \\\bottomrule
\end{tabular}\end{center}

\section{R.79/101 與 R.92/100}
\begin{center}\begin{tabular}{llll}\toprule 韻類 & 藏文 & GX & GHC \\\midrule
R.79 & རྸེ & \emph{-er} & \emph{-jijr} \\ R.101 & འརྸེ & \emph{-er} & \emph{-jiir} \\ R.92 & རྸྀ & \emph{-ər} & \emph{-jɨr} \\ R.100 & འརྸྀ & \emph{-ər} & \emph{-jɨɨr} \\\bottomrule
\end{tabular}\end{center}

GX 本身合併 R.79/R.101 與 R.92/R.100；本方案則以有無前加འ保存傳統韻書中的區別。資料依據：\href{https://homepage.univie.ac.at/xun.gong/tangut/phonology-202409.html}{Gong Xun, Tangut phonology}及\href{https://semakosa.github.io/tangut-pronunciation-db/docs/GX202411-zh.html}{GX202411}。
\end{document}
